\scanpage{146}
% archive foliation: 145
Sur l'ens des \deleted{structures} classes d'isomorphie de structures de Torelli directes, le groupe
$\deleted{\unsure{\mathrm{Sl}(2)}}$
\begin{equation}
\mathrm{Sl}(2,\mathbb Z)=
\left\{\begin{pmatrix}a&b\\c&d\end{pmatrix}
\ \middle|\ 
\begin{gathered}
a,b,c,d\in\mathbb Z\\
ad-bc=1
\end{gathered}\right\}
\tag{9}
\end{equation}
opère à droite, par
\[
t\mathrel{\unsure{\mapstochar\relbar\relbar|\longmapsto}}t\circ g,
\]
on va expliciter cette opération, en tant qu'opération sur $\mathcal D^+$, et $=$ sur $\unsure{\mathcal X|\mathfrak X}_{\mathcal D^+}$ (vu que l'on a une opération de $\mathrm{Sl}(2,\mathbb Z)$ sur la situation modulaire).

On voit qu'à une structure de Torelli $(E,(z_1,z_2))$, $g$ associe la structure $(E,(z'_1,z'_2))$, avec
\[
\begin{aligned}
z'_1&=az_1+cz_2\\
z'_2&=bz_1+dz_2
\end{aligned}
\]
d'où, si $\tau,\tau'$ sont les invariants correspondants
\[
-\qquad\tau=z_2/z_1,\quad\tau'=z'_2/z'_1
\]
la formule
\begin{equation}
\tau'=\frac{d\tau+b}{\unsure{c}\tau+a}
\qquad
(\tau'=\tau\circ g,\ g=\begin{pmatrix}a&b\\c&d\end{pmatrix})
\tag{10}
\end{equation}
Si on remplace l'opération envisagée par l'opération à gauche correspondante, \deleted{\unsure{on}} $t\mathrel{\unsure{\mapstochar\relbar\relbar|\longmapsto}}t\circ g^{-1}$, notée aussi $\tau\longmapsto\tau\circ g^{-1}$, on trouve,
